\documentclass[11pt,a4paper]{article}
\usepackage[T1]{fontenc}
\usepackage{lmodern,microtype,amsmath,amssymb,amsthm,booktabs}
\usepackage[margin=25mm]{geometry}
\usepackage[colorlinks=true,linkcolor=blue,citecolor=blue,urlcolor=blue]{hyperref}
\newtheorem{theorem}{Theorem}
\newtheorem{lemma}[theorem]{Lemma}
\newtheorem{corollary}[theorem]{Corollary}
\newcommand{\E}{\mathbb E}
\newcommand{\Prb}{\mathbb P}
\newcommand{\Law}{\operatorname{Law}}
\newcommand{\cc}{C_n^{\mathrm{causal}}}
\newcommand{\ii}{\imath}
\newcommand{\MEC}{\operatorname{MEC}}
\title{An Exact Causal Construction Below the Logarithmic Scale\\[3pt]
\large Equal entropy, exact dyadic refinement, and a martingale envelope}
\author{A constructive note for the history-indexed C3-C model}
\date{26 September 2026}
\begin{document}
\maketitle
\begin{abstract}
We construct an entropy-optimal finite discrete seed for each horizon $n$, valid under every
adaptive action policy, for an explicit nonpermutation pair with equal
entropy $5$ and equal varentropy $1$. Its entropy is
$5n+O(\sqrt{\log(n+1)})$. The construction uses exact dyadic partitions,
not independent block couplings. Its information distribution is a Bellman
envelope of controlled martingale sums. A uniform martingale Berry--Esseen
bound with vanishing conditional third moments, combined with a tail
estimate, proves the sublogarithmic overhead. Bounded overhead is suggested
by numerical values but is not proved here.
\end{abstract}

\section{The pair and the result}
All entropies and information contents use base-two logarithms. Define
\begin{align*}
P&=\bigl(1/8,\underbrace{1/32,\ldots,1/32}_{24\ \mathrm{symbols}},
              \underbrace{1/128,\ldots,1/128}_{16\ \mathrm{symbols}}\bigr),\\
Q&=\bigl(\underbrace{1/16,\ldots,1/16}_{8\ \mathrm{symbols}},
              \underbrace{1/64,\ldots,1/64}_{32\ \mathrm{symbols}}\bigr).
\end{align*}
Their information contents have laws
\[
\mu_P=\tfrac18\delta_3+\tfrac34\delta_5+\tfrac18\delta_7,
\qquad \mu_Q=\tfrac12\delta_4+\tfrac12\delta_6.
\]
In particular, their entropies are both $5$, their varentropies are both
$1$, and their third centered information moments are both zero.
Their positive support sizes, $41$ and $40$, differ.

\begin{theorem}\label{thm:main}
For this fixed pair there is an explicit sequence of finite source laws
$R_n$ and deterministic response maps, exact under every adaptive policy,
such that
\[
\boxed{\frac12\le \cc(P,Q)-5n
= H(R_n)-5n=O\bigl(\sqrt{\log(n+1)}\bigr).}
\]
Consequently a logarithmic upper bound holds for this nonpermutation,
equal-entropy pair. The construction actually has $o(\log n)$ overhead.
\end{theorem}

Here $\cc$ is the minimum seed entropy in the history-indexed causal
model C3-C: one initial seed determines all responses, and under every
adaptive policy each next output has exactly the law selected by the
current action. The seed law may depend on the horizon. No synchronous
cross-history identity or infinite-horizon consistency is imposed.

\section{An exact partition lemma}
In this note a \emph{dyadic law} means a finite probability vector whose
every positive atom is $2^{-k}$ for a nonnegative integer $k$.
For such a law $A$, let
\[
F_A(k)=\Prb\{\ii_A(X)\le k\},\qquad
\ii_A(x)=-\log_2 A(x),\quad X\sim A.
\]
Say that $R$ refines $A$ if there is a deterministic map taking a random
variable of law $R$ to one of law $A$.

\begin{lemma}[Dyadic refinement criterion]\label{lem:partition}
For finite dyadic laws $R,A$, a refinement map exists if and only if
\[
F_R(k)\le F_A(k)\qquad\hbox{for every integer }k.
\]
It can be constructed by assigning source atoms in decreasing order of
size to target bins with sufficient remaining capacity.
\end{lemma}
\begin{proof}
An atom of size at least $2^{-k}$ must be assigned to a target bin of
size at least $2^{-k}$. This proves necessity.
For sufficiency, process source atoms by increasing information length.
When assigning atoms of mass $2^{-k}$, all remaining capacities of target
bins originally of size at least $2^{-k}$ are integer multiples of
$2^{-k}$. Their combined remaining capacity is
$F_A(k)-F_R(k-1)$, at least the total source mass
$F_R(k)-F_R(k-1)$ now being assigned. Thus every atom fits.
After all source atoms have been assigned, equal total masses imply that
every target bin is exactly filled.
\end{proof}

\begin{lemma}[The minimum CDF is realizable]\label{lem:min}
For two finite dyadic laws $A,B$, the function
$F(k)=\min\{F_A(k),F_B(k)\}$ is the information CDF of a finite dyadic
law $R$ refining both $A$ and $B$.
\end{lemma}
\begin{proof}
The function is nondecreasing, is zero for $k<0$, and eventually equals
$1$. Both $F_A(k)$ and $F_B(k)$ are integer multiples of $2^{-k}$.
Consequently
\[
N_k=2^k\bigl(F(k)-F(k-1)\bigr)
\]
is a nonnegative integer. Take $N_k$ atoms of mass $2^{-k}$ at each
level $k$. Their total mass is one and their information CDF is $F$.
Lemma~\ref{lem:partition} supplies both refinement maps.
\end{proof}

There is no rounding penalty in this step. In fact this law attains the
minimum entropy coupling value:
\[
\MEC(A,B)=\tfrac12\bigl(H(A)+H(B)+W_1(\mu_A,\mu_B)\bigr).
\]
For the lower bound, any common seed atom has probability at most either
of its image probabilities, so its information dominates the maximum of
the two image information contents. The minimum expected maximum is
their common-quantile value, equal to the displayed expression.
The law in Lemma~\ref{lem:min} has precisely this information CDF and
entropy, proving attainability.

\section{A recursive source with exact continuation laws}
Let $R_0$ be a point mass. Recursively define $R_n$ by
\begin{equation}\label{eq:recursion}
F_n(k)=\min\bigl\{(\mu_P*F_{n-1})(k),
                         (\mu_Q*F_{n-1})(k)\bigr\},
\qquad F_0(k)=\mathbf1_{\{k\ge0\}}.
\end{equation}
Convolution here means $(\mu*F)(k)=\sum_\ell\mu(\ell)F(k-\ell)$.
The two CDFs on the right are the information CDFs of
$P\otimes R_{n-1}$ and $Q\otimes R_{n-1}$. Lemma~\ref{lem:min}
constructs $R_n$ and deterministic maps
\begin{equation}\label{eq:decoder}
g_{n,P}:R_n\longrightarrow P\otimes R_{n-1},\qquad
g_{n,Q}:R_n\longrightarrow Q\otimes R_{n-1}.
\end{equation}

Start with one sample of $R_n$. When an action $a$ is selected, apply
$g_{n,a}$, output its first coordinate, and retain its second coordinate
as the new state. Repeat with $n-1$ remaining steps.
Conditional on the output, the retained state has \emph{exactly} law
$R_{n-1}$. Induction shows that, conditional on every positive-probability
history, the current state has the required continuation law. Therefore
the simulator is exact for every adaptive policy. No further random bits
are injected after the initial seed.

\begin{theorem}[Exact optimality of the CDF recursion]\label{thm:optimal}
For any finite pair of dyadic action laws, the construction
\eqref{eq:recursion} satisfies
\[
C_n^{\mathrm{causal}}=H(R_n).
\]
More strongly, every admissible discrete seed $S$ has information CDF
$F_S(t)\le F_n(t)$ for all real $t$. The competing seed need not be dyadic.
\end{theorem}
\begin{proof}
At horizon zero every seed has nonnegative information, so its CDF is at
most $F_0(t)=\mathbf1_{\{t\ge0\}}$. Assume the assertion for horizon
$n-1$, and fix an arbitrary admissible horizon-$n$ seed $S$.
Choose the first action $a$. Its output $Y_a$ is a deterministic function
of $S$. Conditional on each output $y$ of positive probability, the seed
law $S\mid Y_a=y$ is an admissible horizon-$(n-1)$ seed: every continuation
policy must still give the prescribed memoryless conditional output laws.
Since, for each atom in that fiber,
\[
-\log_2\Prb(S=s)
=-\log_2 P_a(y)-\log_2\Prb(S=s\mid Y_a=y),
\]
the induction hypothesis gives
\[
F_S(t)=\sum_yP_a(y)F_{S\mid Y_a=y}\bigl(t-\ii_{P_a}(y)\bigr)
\le(\mu_{P_a}*F_{n-1})(t).
\]
This holds for both actions, hence $F_S(t)\le F_n(t)$.
Integrating survival functions gives $H(S)\ge H(R_n)$, including the
case of infinite seed entropy. The exact decoder construction attains
$H(R_n)$, proving equality.
\end{proof}

Thus the high-dimensional optimization over causal trees reduces, for
dyadic action laws, to the one-dimensional recursion \eqref{eq:recursion}.
This optimality argument uses no normal approximation or moment matching.

For the present pair all calculations can use integers. The CDF of
$R_n$ has denominator $8^n$. On the centered grid $x=-2n,\ldots,2n$,
write $G_n(x)=F_n(5n+x)$. Then
\begin{equation}\label{eq:centered}
G_n(x)=\min\left\{
\frac{G_{n-1}(x+2)+6G_{n-1}(x)+G_{n-1}(x-2)}8,
\frac{G_{n-1}(x+1)+G_{n-1}(x-1)}2\right\}.
\end{equation}
Outside its finite support, each CDF is extended by zero and one.
Its atom counts are
\[
N_{n,k}=2^k\bigl(F_n(k)-F_n(k-1)\bigr).
\]
This is a compact exact description of the seed distribution. Expanding
the atoms and decoder maps can require exponentially many states; no
polynomial-size expanded simulator is claimed.

\section{The martingale envelope and its entropy bound}
Let $\alpha_P$ be the centered increment law taking $-2,0,2$ with
probabilities $1/8,3/4,1/8$, and let $\alpha_Q$ take $-1,1$ with
probabilities $1/2,1/2$. A controller selects either law predictably,
then observes the increment. For every such policy $\pi$, write
$S_n^\pi=D_1+\cdots+D_n$.

\begin{lemma}[Bellman representation]\label{lem:bellman}
For every real $x$,
\begin{equation}\label{eq:envelope}
G_n(x)=\inf_\pi\Prb\{S_n^\pi\le x\}.
\end{equation}
\end{lemma}
\begin{proof}
At horizon zero the assertion is the indicator CDF.
At the first step choose an increment law. Once an increment $d$ is
observed, optimize the remaining probability at threshold $x-d$.
This gives exactly \eqref{eq:centered}. The finite horizon and finite
increment sets allow the conditional minima to be attained.
\end{proof}

The policy attaining the infimum may depend on $x$. We therefore need
a normal approximation uniform over \emph{all} policies, not a
Wasserstein estimate for one policy at a time.
For every policy and every step,
\begin{equation}\label{eq:moments}
\E[D_i\mid\mathcal F_{i-1}]=0,\quad
\E[D_i^2\mid\mathcal F_{i-1}]=1,\quad
\E[D_i^3\mid\mathcal F_{i-1}]=0,\quad |D_i|\le2.
\end{equation}
Apply the martingale Berry--Esseen theorem of Wu, Ma, Sang and Fan
\cite[Theorem 1.1]{wu} to $\xi_i=D_i/\sqrt n$.
The predictable quadratic variation equals $1$ exactly. With their
parameter $\rho=1$,
\[
\E[|\xi_i|^4\mid\mathcal F_{i-1}]
\le\frac4n\E[\xi_i^2\mid\mathcal F_{i-1}],
\]
so their assumptions hold with $\varepsilon_n=2/\sqrt n$ and
$\delta_n=0$ for $n\ge16$. Thus an absolute constant $K$ satisfies
\begin{equation}\label{eq:be}
\sup_\pi\sup_{x\in\mathbb R}
\left|\Prb\{S_n^\pi\le x\}-\Phi(x/\sqrt n)\right|
\le\frac K{\sqrt n}.
\end{equation}
By \eqref{eq:envelope}, the same bound holds with $G_n(x)$ in place
of the policy CDF. This use of a published probability theorem is the
only non-elementary input to the rate estimate.

The bounded-increment exponential estimate gives, for $x>0$, uniformly
over policies,
\[
\Prb\{S_n^\pi>x\},\ \Prb\{S_n^\pi\le-x\}
\le e^{-x^2/(8n)}.
\]
For example, it follows by iterating the conditional estimate
$\E[e^{tD_i}\mid\mathcal F_{i-1}]\le e^{2t^2}$ and optimizing $t$.
Therefore both tails of the distribution with CDF $G_n$ obey the same
bound. If $Z_n$ has this CDF, then $H(R_n)=5n+\E Z_n$.
For any $a>0$, the one-dimensional CDF formula for $W_1$ yields
\begin{align}
0\le\E Z_n
&\le W_1\bigl(\Law(Z_n),\mathcal N(0,n)\bigr)\notag\\
&=\int_{\mathbb R}|G_n(x)-\Phi(x/\sqrt n)|\,dx\notag\\
&\le\frac{2aK}{\sqrt n}
       +4\int_a^\infty e^{-x^2/(8n)}\,dx\notag\\
&\le\frac{2aK}{\sqrt n}+\frac{16n}{a}e^{-a^2/(8n)}.
\label{eq:integrated}
\end{align}
Set $a=\sqrt{8n\ln(n+1)}$. The first term is
$4\sqrt2K\sqrt{\ln(n+1)}$ and the second is bounded. Increasing the
constant for $n<16$ proves the upper bound in Theorem~\ref{thm:main}.
In particular, no interchange of an integral and a policy supremum is
being assumed: the CDF and tail bounds hold uniformly before integration.

\paragraph{The positive lower bound.}
For one step, $P$-information is always odd and $Q$-information is always
even. In any coupling their absolute difference is therefore at least
$1$. The common-seed atom inequality used above gives
$C_1^{\mathrm{causal}}\ge5+1/2$.
The constructed $R_1$ has information masses
$(1/8,3/8,3/8,1/8)$ at lengths $(4,5,6,7)$ and entropy $5+1/2$,
so equality holds. Finally, conditioning any horizon-$(n+1)$ seed on
its first output under action $P$ gives
\[
C_{n+1}^{\mathrm{causal}}\ge5+C_n^{\mathrm{causal}}.
\]
Hence $C_n^{\mathrm{causal}}-5n\ge1/2$ for every $n$.

\section{Scope, exact checks, and the remaining question}
The same argument applies to any finite pair of dyadic laws with equal
entropy, a common positive varentropy, and zero third centered information
moment for each action: bounded support supplies the fourth-moment and
tail estimates, and dyadic refinement supplies exact continuation maps.
The zero-varentropy case with equal entropy consists of uniform laws on
the same number of positive-probability symbols and has zero overhead.

Let $E_n=H(R_n)-5n$ denote the optimal residual, by
Theorem~\ref{thm:optimal}.
\begin{center}
\begin{tabular}{rll}
\toprule
$n$ & $E_n$ (bits) & computation\\
\midrule
1 & $0.5$ & exact\\
2 & $0.5625$ & exact\\
5 & $0.658203125$ & exact\\
10 & $0.689732807688415\ldots$ & exact rational, decimal shown\\
100 & $0.734055223902409\ldots$ & exact rational, decimal shown\\
128 & $0.736278139421860\ldots$ & exact rational, decimal shown\\
5000 & $0.750361245574903\ldots$ & floating-point experiment\\
\bottomrule
\end{tabular}
\end{center}
The script checks exact dyadic histograms through horizon $128$ and
constructs full deterministic decoder maps at horizons $1$ and $2$.
The latter have $54$ and $2312$ seed atoms. Every output/continuation
pair is checked with integer arithmetic to have exactly the prescribed
product probability. This certificate verifies the recursion directly;
the all-horizon claim follows from the partition lemmas and induction.

The numerical values suggest bounded overhead, but this note proves only
$O(\sqrt{\log(n+1)})$. It does not prove a finite limit of $E_n$.
It does prove exact optimality of $R_n$ and the requested logarithmic upper
bound, with a stronger asymptotic estimate, for the displayed equal-entropy
pair. It does not change the square-root obstruction for pairs with equal
entropy and unequal varentropy.

The source is not constrained to be a uniform $5n$-bit block plus a
nonnegative reserve: its individual information lengths lie on both sides
of $5n$. The controlled martingale above is a representation of its
information CDF, not an additional randomness source used by the simulator.
This is why the construction avoids the repeated independent-block costs.
No historical priority claim is made without a separate literature review.

\begin{thebibliography}{9}
\bibitem{wu}
S. Wu, X. Ma, H. Sang, and X. Fan.
\emph{A Berry--Esseen bound of order $1/\sqrt n$ for martingales}.
2020, Theorem 1.1.
\url{https://arxiv.org/abs/2002.00307}.
\end{thebibliography}
\end{document}
